% ==STYLE-BEGIN==
\documentclass[UTF8,11pt]{ctexart}
\usepackage{amsmath,amssymb,amsthm}
\usepackage[a4paper,margin=2.4cm,top=2.3cm,bottom=2.3cm]{geometry}
\usepackage{fancyhdr}
\usepackage{enumitem}
\usepackage{hyperref}
\hypersetup{hidelinks}

% ---- 版式 ----
\linespread{1.15}
\setlength{\parskip}{0.35em}
\setlength{\parindent}{1.6em}
\setlist{nosep,leftmargin=2em}

% ---- 页眉 ----
\pagestyle{fancy}
\fancyhf{}
\fancyhead[L]{\small\coursename}
\fancyhead[R]{\small\thepage}
\renewcommand{\headrulewidth}{0.4pt}

% ---- 定理环境：中文用正体，标题加粗，按小节共用一条流水号 ----
\newtheoremstyle{cnthm}{0.7em}{0.7em}{\normalfont}{}{\bfseries}{}{0.5em}{}
\theoremstyle{cnthm}
\newtheorem{thm}{定理}[section]
\newtheorem{lem}[thm]{引理}
\newtheorem{prop}[thm]{命题}
\newtheorem{cor}[thm]{推论}
\newtheorem{defn}[thm]{定义}
\newtheorem{rem}[thm]{注}
\newtheorem{exa}[thm]{例}
\newtheorem{conj}[thm]{猜想}
\renewcommand{\proofname}{\bfseries 证明}
\providecommand{\coursename}{}
% ==STYLE-END==
\renewcommand{\coursename}{数学分析\,·\,0923}
\title{数学分析 0923}
\date{}
\begin{document}
\maketitle

\section{\ Recall}

原命题 $P\Rightarrow Q$；逆命题 $Q\Rightarrow P$；逆否命题 $\neg Q\Rightarrow\neg P$。

反证法的原理即“原与逆否，同真同假”：通过否定结论 $\neg Q$ 推出条件不成立 $\neg P$，
达到证明原命题的目的。

\section{\ 命题}

\subsubsection*{全称量词 $\forall$}

任意性（对于 $s\in S$，随便取）；一旦取出，就固定。

e.g. $\forall n\in\mathbb Z$ 且 $n\neq0$，$\dfrac{1}{n^2}>0$。

随便取 $n_0\in\mathbb Z$ 且 $n_0\neq0$；一旦取出就固定 $n_0$，接下来只需要验证
$\dfrac{1}{n_0^2}>0$。

\subsubsection*{存在量词 $\exists$}

e.g. $\{$一班学生都擅长数学$\}$ 用数学语言表示。

$x$：某个人；$S=\{$一班学生$\}$；$P(x)$：擅长数学。我们有
\[
\forall x\in S,\quad P(x).
\]
它的否命题是“存在一个一班学生，不擅长数学”，即
\[
\exists x\in S,\quad \neg P(x).
\]

\begin{itemize}
\item $\forall$ 的否命题，要把 $\forall$ 变成 $\exists$；
\item \textbf{否定的不是条件，是结论}。
\end{itemize}

回顾否定命题的方法：

\begin{enumerate}
\item 把 $\forall$ 变成 $\exists$，$\exists$ 变成 $\forall$；
\item 顺序不变；
\item 否定最终命题 $P$。
\end{enumerate}

上面已经对 1、3 两条做出说明，下面来说明第二条。

假设有下面的命题 $\forall x,\forall y,P(x,y)$，现否定此命题。

分析：把后面的 $\forall y,P(x,y)$ 看作是一个命题 $Q(x)$，整个命题变为 $\forall x,Q(x)$。
不难得出，否命题是 $\exists x,\neg Q(x)$。此时单独看 $Q(x)$，有
$\neg Q(x)=\exists y,\neg P(x,y)$。最后整理得否命题是
\[
\exists x,\ \exists y,\ \neg P(x,y).
\]

\section{\ 牛顿二项式展开}

\[
(x+y)^n=\sum_{i=0}^{n}\binom{n}{i}x^iy^{n-i}.
\]

\section{Chapter 1\quad 函数}

\subsubsection*{1.\ 基本定义}

\begin{itemize}
\item \textbf{函数}：一个特殊的映射，$f:X\to\mathbb R$，$X\subseteq\mathbb R$。
\item \textbf{定义域}：使对应法则 $f$ 有意义的点组成的集合。
\item \textbf{函数图像}：
\begin{itemize}
\item 符号：$\mathbb R^2=\{(x,y)\mid x,y\in\mathbb R\}$。
\item 定义：给出 $f:X\to\mathbb R$ 函数，定义 $f$ 的图像为
$\{(x,y)\in\mathbb R^2\mid x\in X,\ y=f(x)\}$。
\end{itemize}
\end{itemize}

一个函数一定是有图像的，如狄利克雷函数，只是我们画不出来罢了。
注意图像与值域不是一回事：图像是平面上的点集，值域是实数的集合。

\subsubsection*{2.\ 奇偶性}

有函数 $f:X\to\mathbb R$，如果定义域 $X$ 关于原点对称，我们可以讨论它的奇偶性。

\begin{itemize}
\item 如果 $\forall x\in X,\ f(x)=f(-x)$，则称 $f(x)$ 是偶函数，
其图像关于 $y$ 轴对称；
\item 如果 $\forall x\in X,\ f(x)=-f(-x)$，则称 $f(x)$ 是奇函数，
其图像关于原点对称。
\end{itemize}

定义里的“$\forall x\in X$”不能省。原笔记写作“如果 $x\in X,\ f(x)=f(-x)$”，
字面上只是一个条件式，不是对定义域中每个点都成立的要求，会变成“存在 $x$ 满足…”，
那就不是偶函数了。

\subsubsection*{3.\ 周期性}

有函数 $f:\mathbb R\to\mathbb R$，若 $\exists T>0$，s.t. $\forall x\in\mathbb R$，
$f(x+T)=f(x)$，则称 $T$ 为 $f$ 的周期。

$\forall x\in\mathbb R,\ \exists T\geq0$，s.t. $f(x+T)=f(x)$ 不是周期性的定义，
这样会使 $T$ 受到 $x$ 的限制。

e.g. $f(x)=\sin x$ 是周期函数，所有的 $2k\pi\ (k\in\mathbb Z^+)$ 都是它的周期。

\textbf{Def:} 如果对于 $f$ 所有的周期，存在一个最小的正数 $T_0$，则称 $T_0$ 为最小正周期。

不是所有周期函数都有最小正周期，比如狄利克雷函数，所有有理数都是它的周期。
显然，它没有最小正周期。

\subsubsection*{4.\ 单调性}

设 $f:X\to\mathbb R$ 为一个函数，如果 $\forall x_1,x_2\in X$，当 $x_1\leq x_2$ 时，
$f(x_1)\leq f(x_2)$，则称 $f$ 在 $X$ 上单调递增；特别地，如果对于所有 $x_1<x_2$，
都有 $f(x_1)<f(x_2)$，则称 $f$ 在 $X$ 上严格单调递增。

单调性要强调是在哪个区间上增或减。

单调递增的等价形式：

\begin{itemize}
\item $\forall x_1<x_2,\ f(x_1)<f(x_2)$（严格单调递增）；
\item $\forall x_1,x_2$，$x_1-x_2$ 与 $f(x_1)-f(x_2)$ 同号；
\item $\forall x_1,x_2$，$\bigl(x_1-x_2\bigr)\bigl(f(x_1)-f(x_2)\bigr)\geq0$；
\item $\forall x_1\neq x_2$，$\dfrac{f(x_1)-f(x_2)}{x_1-x_2}\geq0$。
\end{itemize}

原笔记第 1 条写的是“$\forall x_1\leq x_2,\ f(x_1)\neq f(x_2)$”，那其实是严格单调递增；
现已还原成严格单调递增的写法。第 2、3、4 条只对 $x_1,x_2$ 作要求，
没有限定 $x_1\leq x_2$，靠“同号”和“$\geq0$”把这个信息装进去；
如果把不等号改成 $>0$，得到的才是严格单调递增。

\subsubsection*{5.\ 有界性}

\textbf{Def:} $f:X\to\mathbb R$ 有界是指 $\exists M>0$，s.t. $\forall x\in X$ 有
$|f(x)|\leq M$，则称 $f(x)$ 在 $X$ 上有界。

反过来当然也不行，如 $\forall x\in X,\ \exists M>0$，s.t. $|f(x)|\leq M$ 会使 $M$ 的值
依赖于 $x$，错误。（原笔记这里把式子误写成了 $f(x+T)=f(x)$，那是周期性的式子。）

否命题即为无界的定义：$\forall M>0,\ \exists x\in X$，s.t. $|f(x)|>M$。

\end{document}

\end{document}
